\documentclass[aspectratio=169,11pt]{beamer}

% ---------------------------------------------------------------------------
% Thème très sobre : pas de couleurs vives, pas de symboles de navigation,
% pas d'ombres, une seule teinte de gris pour la structure.
% ---------------------------------------------------------------------------
\usetheme{default}
\usecolortheme{default}

\setbeamertemplate{navigation symbols}{}
\setbeamertemplate{headline}{}
\setbeamertemplate{blocks}[rounded=false,shadow=false]
\setbeamertemplate{itemize items}[circle]
\setbeamertemplate{itemize subitem}[--]
\setbeamertemplate{footline}{%
  \hfill\usebeamerfont{page number in head/foot}%
  \usebeamercolor[fg]{page number in head/foot}%
  \insertframenumber\,/\,\inserttotalframenumber\kern2em\vskip4pt
}

\setbeamercolor{title}{fg=black}
\setbeamercolor{frametitle}{fg=black}
\setbeamercolor{structure}{fg=black!65}
\setbeamercolor{block title}{bg=black!8,fg=black}
\setbeamercolor{block body}{bg=black!3,fg=black}
\setbeamercolor{item}{fg=black!65}
\setbeamerfont{frametitle}{series=\bfseries,size=\large}
\setbeamerfont{title}{series=\bfseries}

\usepackage[french]{babel}
\usepackage[utf8]{inputenc}
\usepackage[T1]{fontenc}
\usepackage{booktabs}
\usepackage{array}
\usepackage{tikz}
\usetikzlibrary{arrows.meta,positioning}
\usepackage{listings}

% Style de code minimal, monochrome — cohérent avec le thème sobre
\lstset{
  basicstyle=\ttfamily\scriptsize,
  breaklines=true,
  frame=single,
  framerule=0.3pt,
  rulecolor=\color{black!40},
  backgroundcolor=\color{black!3},
  showstringspaces=false,
  commentstyle=\itshape\color{black!50},
  keywordstyle=\bfseries,
  columns=fullflexible,
}

\title{Apprentissage fédéré pour l'IoT médical}
\subtitle{Du cadre de gouvernance HEART aux études de cas empiriques}
\author{Clément Frassier}
\institute{Stage --- UKM / UEL \\ Encadrante : Asma Abu Samah}
\date{\today}

\begin{document}

% =============================================================================
\begin{frame}[plain]
  \titlepage
\end{frame}

\begin{frame}{Plan}
  \tableofcontents
\end{frame}

% =============================================================================
\section{Le cadre posé au début du stage : la gouvernance HEART}

\begin{frame}{Pourquoi commencer par la gouvernance ?}
  \begin{itemize}
    \item Avant d'écrire une ligne de code d'apprentissage fédéré : anticiper
      les risques d'un dispositif porté en continu, collectant des données
      physiologiques sensibles.
    \item Document produit en tout début de stage (\texttt{HEART\_PROP}) :
      matrice de risques, typologie de biais spécifique au wearable, pipeline
      de gouvernance de bout en bout.
    \item Sert de fil conducteur à toute la suite : chaque choix technique
      présenté dans cette présentation répond à une case précise de cette
      matrice, pas à un choix arbitraire.
  \end{itemize}
\end{frame}

\begin{frame}{Matrice de risques}
  \centering
  \small
  \begin{tabular}{lcc}
    \toprule
    Risque & Probabilité & Sévérité \\
    \midrule
    Biais physiologique (mélanine, IMC) & Élevée & Critique \\
    Fatigue d'alerte (faux positifs) & Élevée & Majeure \\
    Fuite de données privées & Faible & Critique \\
    Dérive du modèle (\emph{model drift}) & Élevée & Critique \\
    Contraintes TinyML & Moyenne & Moyenne \\
    Manque d'explicabilité & Moyenne & Majeure \\
    Sur-monitoring & Élevée & Majeure \\
    \bottomrule
  \end{tabular}
  \vspace{1em}
  \begin{itemize}
    \small
    \item Sévérité critique même à probabilité faible (fuite de données) :
      la sensibilité des données de santé impose la prudence.
  \end{itemize}
\end{frame}

\begin{frame}{Une typologie de biais propre au wearable}
  \centering
  \small
  \begin{tabular}{lp{5.2cm}}
    \toprule
    Type de biais & Impact sur le projet \\
    \midrule
    Physiologique (mélanine) & Alertes cardiaques moins fiables pour les
      peaux foncées (absorption lumineuse différente) \\
    Composition corporelle & Diagnostic erroné, sensibilité réduite chez
      les patients obèses \\
    Démographique & Modèle optimal pour un cohorte ``standard'', pas pour
      les minorités \\
    Étiquetage & Si les experts cliniques divergent, le modèle apprend une
      ``vérité arbitraire'' \\
    Sélection (pilote) & Performance surestimée : les volontaires ne
      représentent pas la population réelle \\
    \bottomrule
  \end{tabular}
\end{frame}

\begin{frame}{Le pipeline de gouvernance de bout en bout}
  \small
  \begin{description}
    \item[Data] Ownership, minimisation des données de santé, étude
      d'équité, traçabilité des accès.
    \item[Modeling] Architecture \emph{edge-aware}, choix du mode
      d'apprentissage selon la confidentialité (central / fédéré / hybride),
      métriques multidimensionnelles.
    \item[Packaging] Edge vs cloud, conversion pour le matériel cible,
      explicabilité embarquée.
    \item[Validation] Tests bancs, pilotes terrain, cohortes diversifiées.
    \item[Monitoring] Télémétrie, détection de dérive, rounds de
      ré-entraînement fédéré.
    \item[Documentation] Model card, traçabilité, journal de décisions.
  \end{description}
  \vspace{0.5em}
  \begin{block}{Message clé}
    \small La case ``choix du mode d'apprentissage selon la confidentialité''
    est exactement ce qui motive tout le travail empirique qui suit.
  \end{block}
\end{frame}

% =============================================================================
\section{Mon travail --- vue d'ensemble}

\begin{frame}{Ce que j'ai fait}
  \begin{itemize}
    \item Deux études de cas indépendantes d'apprentissage fédéré appliqué à
      des données physiologiques portées : détection du stress chez 15
      infirmières, puis réplication sur un second jeu de données public
      (WESAD) pour tester si les conclusions généralisent.
    \item Pipeline complet reproduit de bout en bout : prétraitement,
      modèle TinyML, client/serveur Flower, confidentialité différentielle,
      personnalisation, grille d'expériences, analyse statistique.
    \item Une investigation méthodologique menée cette semaine sur la
      reproductibilité du bruit DP-SGD, jusqu'à un test de contrôle dédié.
  \end{itemize}
\end{frame}

\begin{frame}{Pourquoi c'est important}
  \begin{itemize}
    \item Répond directement à une case du framework de gouvernance HEART
      posé en début de stage : le choix du mode d'apprentissage selon la
      confidentialité n'est plus théorique, il est mesuré.
    \item Démontre empiriquement --- pas seulement en théorie --- que le FL
      personnalisé peut être compétitif, voire supérieur, au centralisé :
      pas un simple compromis coûteux imposé par la confidentialité.
    \item Des résultats obtenus honnêtement : les erreurs et effets
      inattendus trouvés en cours de route ont été creusés et corrigés,
      pas dissimulés derrière des chiffres qui ``marchent''.
  \end{itemize}
\end{frame}

\begin{frame}{Compétences développées}
  \small
  \begin{description}
    \item[Techniques] Flower~v2, confidentialité différentielle (Opacus),
      contraintes TinyML (quantification, taille mémoire), tests
      statistiques (Welch, Mann--Whitney, correction de Bonferroni).
    \item[Méthodologiques] Diagnostiquer un résultat surprenant jusqu'à un
      mécanisme vérifié plutôt qu'une explication plausible ; concevoir des
      expériences de contrôle pour isoler une cause.
    \item[Outils] Orchestration de jobs longs en arrière-plan, rédaction
      scientifique rigoureuse (LaTeX, gestion de versions).
  \end{description}
\end{frame}

% =============================================================================
\section{Philosophie et mécanismes du Federated Learning}

\begin{frame}{Pourquoi le FL ?}
  \begin{itemize}
    \item Les données physiologiques collectées par bracelet sont des
      données de santé (RGPD, HIPAA).
    \item Les centraliser crée un point de défaillance unique et un risque
      de ré-identification --- exactement la case \emph{``Fuite de données
      privées''} de la matrice HEART.
    \item Le FL inverse le paradigme : le modèle voyage vers la donnée, pas
      l'inverse.
    \item Seules les mises à jour de poids transitent sur le réseau ---
      jamais les signaux physiologiques bruts.
  \end{itemize}
\end{frame}

\begin{frame}{Le cycle FL en quatre étapes}
  \centering
  \begin{tikzpicture}[
      node distance=1.1cm and 1.6cm,
      box/.style={draw=black!60, rounded corners=1pt, align=center,
                   minimum width=3.1cm, minimum height=1.1cm, font=\small},
      arr/.style={-{Latex[length=2mm]}, thick, black!60}
    ]
    \node[box] (srv) {Serveur\\modèle global};
    \node[box, below left=of srv] (c1) {Client 1\\données privées};
    \node[box, below=of srv] (c2) {Client 2\\données privées};
    \node[box, below right=of srv] (c3) {Client 3\\données privées};

    \draw[arr] (srv) -- (c1) node[midway, above left, font=\tiny]{① poids};
    \draw[arr] (srv) -- (c2) node[midway, right, font=\tiny]{① poids};
    \draw[arr] (srv) -- (c3) node[midway, above right, font=\tiny]{① poids};

    \draw[arr] (c1) to[bend right=25] node[midway, below left, font=\tiny]{③ mises à jour} (srv);
    \draw[arr] (c2) to[bend right=15] node[midway, left, font=\tiny]{③} (srv);
    \draw[arr] (c3) to[bend left=25] node[midway, below right, font=\tiny]{③} (srv);
  \end{tikzpicture}
  \vspace{0.5em}
  \small
  ② Entraînement local sur données privées $\quad\to\quad$
  ④ Agrégation $\to$ nouveau modèle global (round suivant)
  \vspace{0.3em}

  \begin{block}{Principe clé}
    \small Les étapes ①→②→③→④ se répètent à chaque round. Le signal brut
    reste sur l'appareil du patient à chaque instant.
  \end{block}
\end{frame}

\begin{frame}{Panorama des mécanismes existants}
  \centering
  \footnotesize
  \begin{tabular}{lll}
    \toprule
    Besoin & Mécanisme & Idée \\
    \midrule
    Hétérogénéité (non-IID) & FedProx & Terme proximal freinant la dérive
      locale \\
    Convergence lente & FedYogi / FedAdam & Optimiseur adaptatif côté
      serveur \\
    Volumes inégaux & FedNova & Normalisation par le nb.\ de pas locaux \\
    Personnalisation & FedPer / Ditto & Tête (ou modèle) locale non
      agrégée \\
    Confidentialité formelle & DP-SGD (Opacus) & Bruit gaussien calibré +
      écrêtage \\
    Clients compromis & FedMedian / TrimmedAvg & Agrégation robuste aux
      valeurs aberrantes \\
    \bottomrule
  \end{tabular}
  \vspace{1em}

  Ce n'est pas une liste exhaustive à mémoriser, mais un arbre de décision :
  chaque contrainte matérielle ou clinique pointe vers un mécanisme précis.
\end{frame}

\begin{frame}{Confidentialité : une défense en couches}
  \small
  \begin{enumerate}
    \item \textbf{Localité des données} --- le FL de base : rien ne quitte
      l'appareil.
    \item \textbf{Sparsification des gradients} --- seules les mises à jour
      les plus significatives transitent.
    \item \textbf{Confidentialité différentielle (DP-SGD)} --- bruit
      calibré, garantie formelle $(\varepsilon,\delta)$.
    \item \textbf{Agrégation sécurisée (SecAgg)} --- le serveur ne voit
      jamais une mise à jour individuelle, seulement la somme agrégée.
  \end{enumerate}
  \vspace{0.5em}
  \begin{block}{Chaque couche est indépendante}
    \small Un objet BLE très contraint peut se limiter aux couches 1--3 ;
    une passerelle smartphone peut se permettre les quatre.
  \end{block}
\end{frame}

\begin{frame}{L'architecture Flower~v2, en bref}
  \small
  \begin{description}
    \item[\texttt{ServerApp}] orchestre les rounds, choisit une stratégie
      (ex.\ FedAvg), évalue le modèle global.
    \item[\texttt{ClientApp}] reçoit les poids, entraîne localement, renvoie
      la mise à jour.
    \item[\texttt{task.py}] toute la logique modèle + données --- \textbf{le
      seul fichier qui change d'un projet à l'autre}.
  \end{description}
  \vspace{1em}
  \centering
  \texttt{flwr run . --stream}
  \vspace{0.5em}

  \footnotesize Une seule commande lance la simulation complète ; le passage
  à des appareils réels ne change que la commande de déploiement, pas le
  code.
\end{frame}

\begin{frame}[fragile]{Hands-on --- le serveur}
\begin{lstlisting}[language=Python]
app = ServerApp()

@app.main()
def main(grid: Grid, context: Context) -> None:
    strategy = FedAvg(fraction_evaluate=1.0)
    result = strategy.start(
        grid=grid,
        initial_arrays=ArrayRecord(Net().state_dict()),
        num_rounds=context.run_config["num-server-rounds"],
        evaluate_fn=global_evaluate,   # eval centralisee apres chaque round
    )
    torch.save(result.arrays.to_torch_state_dict(), "final_model.pt")
\end{lstlisting}
  \footnotesize \texttt{strategy.start()} gère tous les rounds en interne ---
  broadcast, agrégation, évaluation.
\end{frame}

\begin{frame}[fragile]{Hands-on --- le client}
\begin{lstlisting}[language=Python]
app = ClientApp()

@app.train()
def train(msg: Message, context: Context):
    model = Net()
    model.load_state_dict(msg.content["arrays"].to_torch_state_dict())
    pid = context.node_config["partition-id"]
    trainloader, _ = load_data(pid, ...)
    train_fn(model, trainloader, epochs=1,
             lr=msg.content["config"]["lr"])
    return Message(
        content=RecordDict({"arrays": ArrayRecord(model.state_dict())}),
        reply_to=msg)
\end{lstlisting}
  \footnotesize \texttt{partition-id} est injecté automatiquement par
  Flower --- aucune coordination manuelle entre clients.
\end{frame}

\begin{frame}[fragile]{Hands-on --- la configuration}
\begin{lstlisting}
[tool.flwr.app.config]
num-server-rounds = 30
learning-rate     = 0.01
batch-size        = 32

[tool.flwr.federations.local-simulation]
options.num-supernodes = 15
\end{lstlisting}
  \begin{itemize}
    \small
    \item Tout se règle dans \texttt{pyproject.toml}, jamais en dur dans le
      code --- injecté à l'exécution via \texttt{Context}.
    \item Passer de 15 clients simulés à des milliers d'appareils réels ne
      change que ce fichier, pas la logique.
  \end{itemize}
\end{frame}

\begin{frame}[fragile]{Hands-on --- FedProx (hétérogénéité)}
\begin{lstlisting}[language=Python]
global_params = {k: v.clone() for k, v in model.state_dict().items()}
mu = msg.content["config"].get("proximal_mu", 0.01)

for batch in trainloader:
    loss = criterion(model(batch["x"]), batch["y"])
    proximal = sum(torch.sum((p - global_params[n]) ** 2)
                   for n, p in model.named_parameters())
    (loss + (mu / 2.0) * proximal).backward()
    optimizer.step()
\end{lstlisting}
  \footnotesize Le terme proximal pénalise l'écart au modèle global ---
  freine la dérive locale sur des patients hétérogènes.
\end{frame}

\begin{frame}[fragile]{Hands-on --- confidentialité différentielle}
\begin{lstlisting}[language=Python]
pe = PrivacyEngine()   # attache UNE SEULE FOIS par client
net_dp, opt_dp, loader_dp = pe.make_private(
    module=net, optimizer=optimizer, data_loader=trainloader,
    noise_multiplier=0.8,   # controle le compromis confidentialite/precision
    max_grad_norm=1.0,      # ecretage : limite l'influence d'un exemple
)
epsilon = pe.get_epsilon(delta=1e-5)   # garantie cumulative
\end{lstlisting}
  \begin{itemize}
    \small
    \item \texttt{PrivacyEngine} doit être créé une fois, jamais recréé à
      chaque round --- sinon $\varepsilon$ n'est plus cumulé correctement.
    \item Plus \texttt{noise\_multiplier} est élevé, plus la garantie est
      forte, au prix de la précision.
  \end{itemize}
\end{frame}

\begin{frame}[fragile]{Hands-on --- personnalisation (FedPer)}
\begin{lstlisting}[language=Python]
# Envoi au serveur : exclure la tete de classification (fc3)
shared = {k: v for k, v in model.state_dict().items()
          if "fc3" not in k}
model_record = ArrayRecord(shared)
\end{lstlisting}
  \begin{itemize}
    \small
    \item L'extracteur de caractéristiques est partagé et agrégé ; la tête
      de classification reste locale, jamais envoyée.
    \item Permet à chaque patient d'avoir une frontière de décision
      personnalisée, sans réentraîner tout le modèle.
  \end{itemize}
\end{frame}

\begin{frame}{Guide de décision --- hyperparamètres}
  \centering
  \footnotesize
  \begin{tabular}{lll}
    \toprule
    Contrainte & Levier & Valeur recommandée \\
    \midrule
    RAM $<$ 1~Mo & \texttt{batch-size} & 4--16 \\
    Bande passante BLE & sparsité & 0,9--0,995 \\
    $\varepsilon \le 1$ (cible médicale) & \texttt{noise\_multiplier} & $\approx 1{,}8$ (30 rounds) \\
    Patients hétérogènes & \texttt{proximal\_mu} & 0,01--0,05 \\
    Convergence rapide & calendrier de lr & $0{,}01 \to 0{,}001$ après round 5 \\
    \bottomrule
  \end{tabular}
\end{frame}

\begin{frame}{Combinaisons recommandées --- TinyML médical}
  \centering
  \footnotesize
  \begin{tabular}{lp{6cm}}
    \toprule
    Contrainte & Combinaison \\
    \midrule
    RAM $<$ 512~Ko, BLE & FedProx+GroupNorm + DP léger + sparsification
      agressive + FedMedian \\
    RAM 2--8~Mo, WiFi & FedProx+GroupNorm + DP + SecAgg + FedYogi \\
    Passerelle smartphone & Ditto+GroupNorm + DP strict + SecAgg \\
    \bottomrule
  \end{tabular}
  \vspace{1em}

  \footnotesize BatchNorm est incompatible avec Opacus (statistiques
  inter-échantillons) --- GroupNorm est la norme compatible DP recommandée.
\end{frame}

% =============================================================================
\section{Reproduire la méthode, étape par étape}

\begin{frame}{Un pipeline générique en cinq étapes}
  \begin{enumerate}
    \item Partitionner les données par client (1 client = 1 personne).
    \item Définir un modèle TinyML (quelques milliers de paramètres,
      quelques Ko).
    \item Implémenter client/serveur Flower --- code générique, réutilisable.
    \item Construire une grille d'expériences (stratégie $\times$ niveau de
      bruit $\times$ graine).
    \item Évaluer statistiquement --- tests appariés, correction pour
      comparaisons multiples, jamais une seule moyenne isolée.
  \end{enumerate}
\end{frame}

\begin{frame}{Ce qui se réutilise à l'identique}
  \begin{itemize}
    \item \texttt{client\_app.py} et \texttt{server\_app.py} sont
      \textbf{strictement identiques} entre les deux études de cas
      présentées ensuite (infirmières et WESAD).
    \item Seul \texttt{task.py} change (chargement des données + modèle).
    \item C'est la preuve que la méthode n'est pas ajustée \emph{ad hoc}
      pour un seul jeu de données --- elle transfère.
  \end{itemize}
\end{frame}

\begin{frame}{La recette condensée}
  \small
  Pour reproduire cette démarche sur un nouveau jeu de données :
  \begin{enumerate}
    \item Identifier l'unité ``client'' naturelle (patient, sujet, appareil).
    \item Vérifier l'équilibre de classes par client --- un client
      mono-classe fausse toute métrique groupée.
    \item Fixer un modèle minimal compatible microcontrôleur.
    \item Lancer une grille multi-graines, jamais une seule exécution.
    \item Toujours comparer une métrique groupée \emph{et} une métrique
      restreinte aux clients réellement informatifs.
    \item Vérifier la reproductibilité d'une graine avant de lui attribuer
      un comportement fixe --- surtout sous confidentialité différentielle.
  \end{enumerate}
\end{frame}

% =============================================================================
\section{Aperçu des deux études de cas}

\begin{frame}{Étude de cas 1 --- Stress des infirmières}
  \small
  \begin{itemize}
    \item 15 infirmières, bracelet connecté, détection binaire du stress ;
      FedProx+DP puis personnalisation (FedPer, puis tête multi-couches +
      SCAFFOLD).
    \item \textbf{Résultat clé} : l'hétérogénéité extrême entre infirmières
      limite la généralisation de \emph{toute} architecture, centralisée
      comprise (son AUC tombe même sous le niveau du hasard sur les cas
      difficiles). La personnalisation réduit cet écart de façon
      significative, sans l'éliminer.
    \item Le FL en lui-même n'ajoute qu'un désavantage marginal par-dessus
      ce problème de fond --- pas la cause principale.
  \end{itemize}
\end{frame}

\begin{frame}{Étude de cas 2 --- WESAD (réplication indépendante)}
  \small
  \begin{itemize}
    \item Second jeu de données, capteurs poignet uniquement --- teste si
      les conclusions du premier projet généralisent.
    \item \textbf{Résultat clé} : le coût du FL face au centralisé n'est pas
      détectable statistiquement, et la personnalisation (FedPer) ne
      réduit pas seulement l'écart --- elle \textbf{dépasse} le centralisé.
    \item Une investigation cette semaine sur un effet inattendu (jour-1
      instable selon la graine) a mené à un mécanisme confirmé par un test
      de contrôle dédié : le bruit de confidentialité différentielle,
      pas un bug.
  \end{itemize}
\end{frame}

% =============================================================================
\section{Leçons apprises}

\begin{frame}{Leçons positives}
  \begin{itemize}
    \item La méthode transfère telle quelle d'un jeu de données à l'autre
      --- même code client/serveur, seul \texttt{task.py} change.
    \item La personnalisation légère (une seule couche non agrégée) peut
      \textbf{dépasser}, pas seulement approcher, le centralisé.
    \item Le framework de gouvernance posé en tout début de stage a
      effectivement structuré les choix techniques faits ensuite ---
      pas un exercice isolé.
  \end{itemize}
\end{frame}

\begin{frame}{Leçons négatives / difficultés}
  \begin{itemize}
    \item DP-SGD introduit une instabilité largement sous-estimée au
      départ (bimodalité, effondrement de mode).
    \item Une métrique groupée sur des clients mono-classe peut
      \textbf{inventer} un écart dans le mauvais sens, pas seulement le
      masquer.
    \item La reproductibilité d'une ``graine'' n'est pas garantie sous
      confidentialité différentielle --- un piège qu'il faut chercher
      activement, jamais supposer absent.
  \end{itemize}
\end{frame}

\begin{frame}{Ce que je ferais différemment}
  \begin{itemize}
    \item Tester la reproductibilité dès la première grille d'expériences,
      pas seulement quand un résultat surprend.
    \item Prévoir un test de contrôle (ex.\ sans DP) en parallèle de chaque
      nouveau mécanisme introduit, pas après-coup.
    \item Documenter systématiquement les biais démographiques potentiels
      dès le choix du jeu de données --- non testés ici faute de métadonnées
      disponibles.
  \end{itemize}
\end{frame}

% =============================================================================
\section{Conclusion}

\begin{frame}{Synthèse}
  \begin{itemize}
    \item Deux études de cas indépendantes, deux histoires différentes,
      une méthode commune.
    \item \textbf{Infirmières} : l'hétérogénéité extrême limite toute
      architecture y compris centralisée ; la personnalisation réduit
      l'écart sans l'éliminer.
    \item \textbf{WESAD} : le coût du FL n'est pas détectable
      statistiquement ; la personnalisation dépasse le centralisé.
    \item Dans les deux cas : le FL reste la seule approche pratiquement
      viable pour ce type de données de santé, avec un coût quantifié
      honnêtement plutôt que supposé.
  \end{itemize}
\end{frame}

\begin{frame}{Retour à la matrice de gouvernance HEART}
  \small
  \begin{description}
    \item[Adressées empiriquement] Model drift (enquête DP-SGD), TinyML
      constraints (modèle 5,32~Ko), choix du mode d'apprentissage selon la
      confidentialité.
    \item[Encore ouvertes] Biais démographique (pas de métadonnées dans ces
      jeux de données), pilotes terrain, explicabilité embarquée
      (SHAP-lite / LIME jamais testés en pratique).
  \end{description}
\end{frame}

\begin{frame}[plain]
  \centering
  \Large Merci --- Questions
\end{frame}

\end{document}
